% Options for packages loaded elsewhere
\PassOptionsToPackage{unicode}{hyperref}
\PassOptionsToPackage{hyphens}{url}
\PassOptionsToPackage{dvipsnames,svgnames,x11names}{xcolor}
%
\documentclass[
  11pt,
]{article}
\usepackage{amsmath,amssymb}
\usepackage{iftex}
\ifPDFTeX
  \usepackage[T1]{fontenc}
  \usepackage[utf8]{inputenc}
  \usepackage{textcomp} % provide euro and other symbols
\else % if luatex or xetex
  \usepackage{unicode-math} % this also loads fontspec
  \defaultfontfeatures{Scale=MatchLowercase}
  \defaultfontfeatures[\rmfamily]{Ligatures=TeX,Scale=1}
\fi
\usepackage{lmodern}
\ifPDFTeX\else
  % xetex/luatex font selection
\fi
% Use upquote if available, for straight quotes in verbatim environments
\IfFileExists{upquote.sty}{\usepackage{upquote}}{}
\IfFileExists{microtype.sty}{% use microtype if available
  \usepackage[]{microtype}
  \UseMicrotypeSet[protrusion]{basicmath} % disable protrusion for tt fonts
}{}
\makeatletter
\@ifundefined{KOMAClassName}{% if non-KOMA class
  \IfFileExists{parskip.sty}{%
    \usepackage{parskip}
  }{% else
    \setlength{\parindent}{0pt}
    \setlength{\parskip}{6pt plus 2pt minus 1pt}}
}{% if KOMA class
  \KOMAoptions{parskip=half}}
\makeatother
\usepackage{xcolor}
\usepackage[margin=1in]{geometry}
\usepackage{color}
\usepackage{fancyvrb}
\newcommand{\VerbBar}{|}
\newcommand{\VERB}{\Verb[commandchars=\\\{\}]}
\DefineVerbatimEnvironment{Highlighting}{Verbatim}{commandchars=\\\{\}}
% Add ',fontsize=\small' for more characters per line
\newenvironment{Shaded}{}{}
\newcommand{\AlertTok}[1]{\textcolor[rgb]{1.00,0.00,0.00}{\textbf{#1}}}
\newcommand{\AnnotationTok}[1]{\textcolor[rgb]{0.38,0.63,0.69}{\textbf{\textit{#1}}}}
\newcommand{\AttributeTok}[1]{\textcolor[rgb]{0.49,0.56,0.16}{#1}}
\newcommand{\BaseNTok}[1]{\textcolor[rgb]{0.25,0.63,0.44}{#1}}
\newcommand{\BuiltInTok}[1]{\textcolor[rgb]{0.00,0.50,0.00}{#1}}
\newcommand{\CharTok}[1]{\textcolor[rgb]{0.25,0.44,0.63}{#1}}
\newcommand{\CommentTok}[1]{\textcolor[rgb]{0.38,0.63,0.69}{\textit{#1}}}
\newcommand{\CommentVarTok}[1]{\textcolor[rgb]{0.38,0.63,0.69}{\textbf{\textit{#1}}}}
\newcommand{\ConstantTok}[1]{\textcolor[rgb]{0.53,0.00,0.00}{#1}}
\newcommand{\ControlFlowTok}[1]{\textcolor[rgb]{0.00,0.44,0.13}{\textbf{#1}}}
\newcommand{\DataTypeTok}[1]{\textcolor[rgb]{0.56,0.13,0.00}{#1}}
\newcommand{\DecValTok}[1]{\textcolor[rgb]{0.25,0.63,0.44}{#1}}
\newcommand{\DocumentationTok}[1]{\textcolor[rgb]{0.73,0.13,0.13}{\textit{#1}}}
\newcommand{\ErrorTok}[1]{\textcolor[rgb]{1.00,0.00,0.00}{\textbf{#1}}}
\newcommand{\ExtensionTok}[1]{#1}
\newcommand{\FloatTok}[1]{\textcolor[rgb]{0.25,0.63,0.44}{#1}}
\newcommand{\FunctionTok}[1]{\textcolor[rgb]{0.02,0.16,0.49}{#1}}
\newcommand{\ImportTok}[1]{\textcolor[rgb]{0.00,0.50,0.00}{\textbf{#1}}}
\newcommand{\InformationTok}[1]{\textcolor[rgb]{0.38,0.63,0.69}{\textbf{\textit{#1}}}}
\newcommand{\KeywordTok}[1]{\textcolor[rgb]{0.00,0.44,0.13}{\textbf{#1}}}
\newcommand{\NormalTok}[1]{#1}
\newcommand{\OperatorTok}[1]{\textcolor[rgb]{0.40,0.40,0.40}{#1}}
\newcommand{\OtherTok}[1]{\textcolor[rgb]{0.00,0.44,0.13}{#1}}
\newcommand{\PreprocessorTok}[1]{\textcolor[rgb]{0.74,0.48,0.00}{#1}}
\newcommand{\RegionMarkerTok}[1]{#1}
\newcommand{\SpecialCharTok}[1]{\textcolor[rgb]{0.25,0.44,0.63}{#1}}
\newcommand{\SpecialStringTok}[1]{\textcolor[rgb]{0.73,0.40,0.53}{#1}}
\newcommand{\StringTok}[1]{\textcolor[rgb]{0.25,0.44,0.63}{#1}}
\newcommand{\VariableTok}[1]{\textcolor[rgb]{0.10,0.09,0.49}{#1}}
\newcommand{\VerbatimStringTok}[1]{\textcolor[rgb]{0.25,0.44,0.63}{#1}}
\newcommand{\WarningTok}[1]{\textcolor[rgb]{0.38,0.63,0.69}{\textbf{\textit{#1}}}}
\usepackage{longtable,booktabs,array}
\usepackage{calc} % for calculating minipage widths
% Correct order of tables after \paragraph or \subparagraph
\usepackage{etoolbox}
\makeatletter
\patchcmd\longtable{\par}{\if@noskipsec\mbox{}\fi\par}{}{}
\makeatother
% Allow footnotes in longtable head/foot
\IfFileExists{footnotehyper.sty}{\usepackage{footnotehyper}}{\usepackage{footnote}}
\makesavenoteenv{longtable}
\setlength{\emergencystretch}{3em} % prevent overfull lines
\providecommand{\tightlist}{%
  \setlength{\itemsep}{0pt}\setlength{\parskip}{0pt}}
\setcounter{secnumdepth}{-\maxdimen} % remove section numbering
\ifLuaTeX
  \usepackage{selnolig}  % disable illegal ligatures
\fi
\IfFileExists{bookmark.sty}{\usepackage{bookmark}}{\usepackage{hyperref}}
\IfFileExists{xurl.sty}{\usepackage{xurl}}{} % add URL line breaks if available
\urlstyle{same}
\hypersetup{
  pdfauthor={Stefano Alimonti --- stefalimonti@gmail.com},
  colorlinks=true,
  linkcolor={blue},
  filecolor={Maroon},
  citecolor={Blue},
  urlcolor={blue},
  pdfcreator={LaTeX via pandoc}}

\author{Stefano Alimonti --- stefalimonti@gmail.com}
\date{March 2026}

\begin{document}

{
\hypersetup{linkcolor=black}
\setcounter{tocdepth}{2}
\tableofcontents
}
\section{Frobenius Eigenvalues and Gauss Sums from Witt Vector Carry
Arithmetic}\label{frobenius-eigenvalues-and-gauss-sums-from-witt-vector-carry-arithmetic}

\textbf{Author:} Stefano Alimonti \textbf{Affiliation:} Independent Researcher
\textbf{Date:} March 2026

\begin{center}\rule{0.5\linewidth}{0.5pt}\end{center}

\newpage

\subsection{Abstract}\label{abstract}

While the Riemann Hypothesis is proved for function fields \(\mathbb{F}_q[x]\),
where addition is carry-free, the extension to \(\mathbb{Z}\) is obstructed by
carry propagation. By lifting carry arithmetic to the ring of Witt vectors
\(W_n(\mathbb{F}_p)\), we construct a transfer operator whose spectrum encodes
Frobenius eigenvalues of algebraic varieties --- connecting the most elementary
operation in positional arithmetic to the deepest structures of algebraic
geometry.

\textbf{Theorem A.} For \(p = 3\), the characteristic polynomial of the
\(9 \times 9\) multiplicative Witt carry matrix contains the exact Frobenius
polynomial \(\mu^2 + 3\mu + 3\) of the supersingular elliptic curve
\(E: y^2 = x^3 + 2x + 1\) over \(\mathbb{F}_3\), with eigenvalues at the Weil
bound \(\lvert\mu\rvert = \sqrt{3}\).

\textbf{Theorem B.} For \(p = 7\), exact computation over
\(\mathbb{Z}[\omega_7]\) yields a genus-3 Frobenius factor
\((\mu^2 - 7)^2(\mu^2 + 7)\) with 6 eigenvalues at
\(\lvert\mu\rvert = \sqrt{7}\). The supersingular component \(\mu^2 + 7\) is the
minimal polynomial of the quadratic Gauss sum \(g(\chi^3) = i\sqrt{7}\) of the
Legendre symbol, with roots \(\mu = \pm g(\chi^3)\) --- an exact algebraic
identification, not a numerical coincidence.

The carry exponent \(c_2\) over \(\mathbb{F}_p^4\) is perfectly equidistributed
on \(\mathbb{F}_p^{\ast}\), and the Weil-bound eigenvectors are orthogonal to
the all-ones vector: the Frobenius is invisible in the total exponential sum but
emerges in the periodic orbit structure \(\mathrm{Tr}(M^k)\). Dirichlet
character twists recover missing Frobenius eigenvalues at every prime tested
(\(p = 3, 5, 7, 11, 13, 17, 19, 23\)): additive Legendre at \(p = 5\); order-5
at \(p = 11\); order-3 at \(p = 13\). The productive character order varies with
\(p\) in a currently unexplained way --- no simple arithmetic condition predicts
which character is needed.

We conjecture that the carry operator computes (a subfactor of) the Frobenius on
the Artin-Schreier variety \(z^p - z = c_2(x_0, x_1, y_0, y_1)\), where \(c_2\)
is the level-2 Witt multiplication carry dominated by the Fermat quotient
\((x^p y^p - (xy)^p)/p\).

\textbf{Keywords:} Witt vectors, carry arithmetic, Frobenius endomorphism, Gauss
sums, Artin-Schreier varieties, Weil conjectures, transfer operators, Fermat
curves.

\textbf{MSC 2020:} 11S40 (Witt vectors), 11G20 (Curves over finite fields),
11T23 (Exponential sums), 11M38 (Zeta and \(L\)-functions in characteristic
\(p\)).

\begin{center}\rule{0.5\linewidth}{0.5pt}\end{center}

\newpage

\subsection{1. Introduction}\label{introduction}

\subsubsection{1.1 Motivation: Carries as the Obstruction --- and the
Solution}\label{motivation-carries-as-the-obstruction-and-the-solution}

The Riemann Hypothesis for function fields \(\mathbb{F}_q[x]\) was proved by
Weil {[}Wei48{]} using the theory of the Frobenius endomorphism on algebraic
curves. A natural question, raised independently by several authors, is whether
an analogous approach could work for \(\mathbb{Z}\).

The fundamental obstruction is arithmetic: in \(\mathbb{F}_q[x]\), digit-wise
addition and multiplication produce no carries --- coefficients operate
independently in each degree. In \(\mathbb{Z}\), carries propagate information
between digits, creating long-range correlations that have no polynomial
analogue.

This paper shows that carries are not merely an obstruction but contain, in a
precise algebraic sense, the Frobenius eigenvalues themselves. At \(p = 3\), the
carry operator reproduces the Frobenius of a supersingular elliptic curve. At
\(p = 7\), it produces a genus-3 Frobenius factor whose supersingular
eigenvalues are exactly the quadratic Gauss sum of the Legendre symbol. At
\(p = 5\) and \(p = 13\), Dirichlet character twists of the carry operator
recover Frobenius eigenvalues that the untwisted operator misses. The natural
home for this correspondence is the ring of Witt vectors \(W(\mathbb{F}_p)\),
where carries arise from the ghost map inversion and the Frobenius acts as a
shift on ghost components.

\subsubsection{1.2 Context: Carry Operators}\label{context-carry-operators}

The spectral theory of carries was initiated by Holte {[}Hol97{]} and
Diaconis--Fulman {[}DF09{]}, who showed that carry propagation in base-\(b\)
addition defines a Markov chain with eigenvalues \(1/b^k\) for \(k \geq 0\).
These eigenvalues are real and positive --- they capture the contractile
dynamics of carry decay but contain no phase information.

The Frobenius eigenvalues of varieties over finite fields, by contrast, are
complex numbers whose moduli satisfy the Weil bound
\(\lvert\alpha\rvert = q^{w/2}\) (weight \(w\)) and whose phases encode the
deepest arithmetic of the variety. The passage from real carry eigenvalues to
complex Frobenius eigenvalues is the central challenge.

\subsubsection{1.3 Main Results}\label{main-results}

We resolve this challenge by constructing carry operators that live not on the
integers but on the Witt vectors. The key innovation is to weight carry
transitions by the additive character \(\omega = e^{2\pi i/p}\), which converts
the real-valued carry into a complex-valued operator whose spectrum encodes both
moduli and phases.

Status labels are used uniformly as \textbf{Proved}, \textbf{Conditional}, and
\textbf{Open target}.

\begin{longtable}[]{@{}
  >{\raggedright\arraybackslash}p{(\columnwidth - 4\tabcolsep) * \real{0.3333}}
  >{\raggedright\arraybackslash}p{(\columnwidth - 4\tabcolsep) * \real{0.3333}}
  >{\raggedright\arraybackslash}p{(\columnwidth - 4\tabcolsep) * \real{0.3333}}@{}}
\toprule\noalign{}
\begin{minipage}[b]{\linewidth}\raggedright
Result
\end{minipage} & \begin{minipage}[b]{\linewidth}\raggedright
Status
\end{minipage} & \begin{minipage}[b]{\linewidth}\raggedright
Method
\end{minipage} \\
\midrule\noalign{}
\endhead
\bottomrule\noalign{}
\endlastfoot
Theorem A (\(p = 3\) Frobenius) & \textbf{Proved} & Computer-assisted, exact
arithmetic over \(\mathbb{Z}[\omega]\) (Faddeev--LeVerrier + SymPy
verification) \\
Theorem B (\(p = 7\) Frobenius + Gauss sum) & \textbf{Proved} &
Computer-assisted, exact arithmetic over \(\mathbb{Z}[\omega_7]\) \\
Theorem 3.5 (additive negative result) & \textbf{Proved} & Same method \\
Corollary C (Weil bound + point counts) & \textbf{Proved} & Follows from Theorem
A \\
Conjecture 1 (Artin-Schreier Carry) & \textbf{Open} & Supported by Theorems A, B
and \(p = 5\) character twist data \\
Conjecture 2 (Universality via Fermat) & \textbf{Open} & Structural evidence at
\(p = 3, 5, 7\) \\
\end{longtable}

\textbf{Theorem A} (Frobenius at \(p = 3\); computer-assisted, exact arithmetic
over \(\mathbb{Z}[\omega]\)). Let \(p = 3\), \(\omega = e^{2\pi i/3}\), and
define the \(9 \times 9\) matrix \(M\) over \(\mathbb{Z}[\omega]\) by

\[M_{(x_1,y_1),(x_0,y_0)} = \omega^{c_2(x_0, y_0, x_1, y_1)}\]

where the indices range over \(\mathbb{F}_3 \times \mathbb{F}_3\) and
\(c_2(x_0, y_0, x_1, y_1)\) is the level-2 multiplicative Witt carry (Definition
2.3). Then:

\[\chi_M(\mu) = \mu^4 \cdot (\mu^2 + 3\mu + 3) \cdot (\mu^3 - 6\mu^2 + 12\mu - 45)\]

The quadratic factor \(\mu^2 + 3\mu + 3\) is the characteristic polynomial of
Frobenius on the supersingular elliptic curve \(E: y^2 = x^3 + 2x + 1\) over
\(\mathbb{F}_3\).

\textbf{Theorem B} (Frobenius and Gauss Sums at \(p = 7\); computer-assisted,
exact arithmetic over \(\mathbb{Z}[\omega_7]\)). The \(49 \times 49\) carry
matrix at \(p = 7\) has 6 eigenvalues at the Weil bound
\(\lvert\mu\rvert = \sqrt{7}\), and its characteristic polynomial contains the
exact factor

\[(\mu^2 - 7)^2 \cdot (\mu^2 + 7)\]

The factor \(\mu^2 + 7\) has roots \(\pm i\sqrt{7}\), which are exactly the
quadratic Gauss sum \(g(\chi^3)\) of the Legendre symbol
\(\left(\frac{\cdot}{7}\right)\). The factor \((\mu^2 - 7)^2\) captures the
superspecial Frobenius eigenvalues \(\pm\sqrt{7}\), each with multiplicity 2.

\textbf{Corollary C.} The carry operator \(T = M/p^2\) has eigenvalues
satisfying the Weil bound \(\lvert\lambda\rvert = p^{-3/2}\). For \(p = 3\), the
Frobenius pair \(\alpha, \bar{\alpha}\) recovers the point counts
\(N_E(\mathbb{F}_{3^k}) = 3^k + 1 - (\alpha^k + \bar{\alpha}^k)\) via the
Lefschetz trace formula.

\subsubsection{1.4 Significance}\label{significance}

These results demonstrate exact spectral factor matches between carry operators
and Frobenius polynomials:

\begin{enumerate}
\def\labelenumi{\arabic{enumi}.}
\item
  \textbf{Frobenius polynomials appear as exact spectral factors of carry
  matrices.} The characteristic polynomials of Frobenius on certain algebraic
  curves over \(\mathbb{F}_p\) divide the characteristic polynomials of the
  corresponding Witt carry matrices --- an exact algebraic match, verified by
  computer-assisted computation with exact arithmetic. This holds at \(p = 3\)
  (Theorem A), \(p = 7\) (exact genus-3 factor), and via character twists at
  \(p = 5\) and \(p = 13\). Whether this spectral match reflects a deeper
  cohomological relationship (Conjecture 1) remains open.
\item
  \textbf{The multiplicative structure is essential.} We show (Theorem 3.5) that
  the additive Witt carry operator has a different characteristic polynomial
  (\(\mu^6(\mu^3 - 6\mu^2 + 9)\)) with no elliptic curve factor. The Frobenius
  emerges only from multiplication, consistent with its nature as a ring
  endomorphism.
\item
  The characteristic polynomial is defined over \(\mathbb{Q}\). Despite the
  matrix \(M\) living in \(\mathbb{Z}[\omega]\), all coefficients of \(\chi_M\)
  are rational integers (Proposition 3.2, verified also at \(p = 7\)). This
  reveals a Galois symmetry: the spectrum is invariant under
  \(\omega \mapsto \omega^k\) for \(\gcd(k, p) = 1\).
\item
  \textbf{Structural consistency with an Artin-Schreier interpretation
  (Conjecture 1).} The carry exponent \(c_2\) is perfectly equidistributed on
  \(\mathbb{F}_p^{\ast}\) (§5.7), and the Weil-bound eigenvectors are orthogonal
  to the all-ones vector (§5.7). The Frobenius information is hidden in the
  periodic orbit structure of the carry, not in the naive exponential sum ---
  consistent with (but not yet proved to be) the cohomological Frobenius on the
  Artin-Schreier variety \(z^p - z = c_2(\boldsymbol{x})\).
\item
  \textbf{Supersingular eigenvalues are Gauss sums.} At \(p = 7\), the
  eigenvalues \(\pm i\sqrt{7}\) from the factor \(\mu^2 + 7\) are exactly the
  quadratic Gauss sum
  \(g(\chi^3) = \sum_{x=1}^{6} \left(\frac{x}{7}\right) \omega^x = i\sqrt{7}\)
  (§5.6). This is the first identification of a carry eigenvalue with a named
  number-theoretic quantity.
\item
  \textbf{The Fermat curve genus predicts the degree of the Frobenius factor.}
  The level-2 carry is dominated by the Fermat quotient
  \(Q_p(x,y) = (x^p y^p - (xy)^p)/p\). The genus \(g = (p-1)(p-2)/2\) of the
  associated Fermat curve determines the degree \(2g\) of the full Frobenius
  polynomial: for \(p = 3\), \(g = 1\) and the factor is quadratic; for
  \(p = 7\), the \(49 \times 49\) matrix captures 6 of the expected \(2g = 30\)
  eigenvalues.
\item
  \textbf{Base-2 bridge compatibility.} Independent Step-3 operator results in
  the binary D-odd program (carry/Witt intertwiner + exact character-resolvent
  channel) show that the same ``selected spectral factor'' architecture survives
  at \(p=2\) in real form (dominant mode \(1/2\)) while \(p \ge 3\) carries
  phase-bearing Frobenius pairs. This supports a unified viewpoint: arithmetic
  constants are encoded by character-projected resolvent factors, with \(p=2\)
  and \(p\ge3\) as two faces of the same operator mechanism.
\end{enumerate}

\subsubsection{1.5 Outline}\label{outline}

Section 2 recalls Witt vector arithmetic and defines the carry operator,
explaining the role of the additive character as a discrete Fourier transform on
the carry space. Section 3 proves Theorem A via the Faddeev--LeVerrier algorithm
with exact arithmetic over \(\mathbb{Z}[\omega]\). Section 4 identifies the
elliptic curve and verifies the Lefschetz trace formula. Section 5 presents the
full computational landscape: the \(p=2\) golden ratio and Ihara zeta (§5.1),
the cross-prime summary table (§5.2), the exact \(p=7\) factorization and
Theorem B (§5.3), the Dirichlet character twist that rescues \(p=5\) (§5.4), the
level-3/4 rank-collapse negative result (§5.5), the Gauss sum identification at
\(p=7\) (§5.6), the equidistribution and orthogonality structure of \(c_2\)
(§5.7), and the extended twist scan to \(p=11, 13\) (§5.8). Section 6 formulates
the Artin-Schreier/Fermat curve conjecture, explaining the genus obstruction and
the character decomposition approach. Section 7 discusses the overlap transfer
matrix (negative result) and Ruelle--Perron--Frobenius theory. Section 8
presents further conjectures, the cross-prime obstruction to an Euler product,
and open problems.

\newpage

\subsubsection{1.6 Conceptual Architecture}\label{conceptual-architecture}

The following diagram summarizes the flow from elementary carry arithmetic to
the Frobenius eigenvalues, and the three disciplines it bridges:

\small

\begin{verbatim}
                    ┌───────────────────────────────────┐
                    │     POSITIONAL ARITHMETIC         │
                    │   (carries in base-p addition     │
                    │    and multiplication)            │
                    └───────────────┬───────────────────┘
                                    │
                          LIFT TO WITT VECTORS
                         W_n(𝔽_p): ghost map
                        inversion ⟹ carries
                                    │
                    ┌───────────────▼───────────────────┐
                    │   WITT CARRY TRANSFER OPERATOR    │
                    │                                   │
                    │   M_{(x₁,y₁),(x₀,y₀)} = ω^{c₂}    │
                    │                                   │
                    │   ω^{c₂} = Fourier transform      │
                    │   on carry space 𝔽_p              │
                    └──────┬────────────────┬───────────┘
                           │                │
              ┌────────────▼───┐     ┌──────▼─────────────┐
              │  COMBINATORICS │     │ ALGEBRAIC GEOMETRY │
              │                │     │                    │
              │  Diaconis–     │     │  Artin-Schreier    │
              │  Fulman        │     │  variety           │
              │  eigenvalues   │     │  z^p−z = c₂(x,y)   │
              │  |λ| = 1/p^k   │     │  |α| = √p (Weil)   │
              │  (real, §3.5)  │     │  (complex, Thm A/B)│
              └────────┬───────┘     └────────┬───────────┘
                       │                      │
                       │    FERMAT QUOTIENT   │
                       │  (x^p·y^p-(xy)^p)/p  │
                       │         │            │
                       └─────────┼────────────┘
                                 │
                    ┌────────────▼──────────────────────┐
                    │       NUMBER THEORY               │
                    │                                   │
                    │   Gauss sums g(χ^k) and           │
                    │   Jacobi sums J(χ_a, χ_b)         │
                    │   = Frobenius eigenvalues         │
                    │   = spectral components of carry  │
                    │                                   │
                    │   p=3: g=1 (elliptic, Thm A)      │
                    │   p=7: g(χ³)=i√7 (Gauss, Thm B)   │
                    │   p=5: rescued by χ₂ twist        │
                    └───────────────────────────────────┘
\end{verbatim}

\normalsize

\begin{center}\rule{0.5\linewidth}{0.5pt}\end{center}

\newpage

\subsection{2. Witt Vector Carry Arithmetic}\label{witt-vector-carry-arithmetic}

\subsubsection{2.1 Truncated Witt Vectors}\label{truncated-witt-vectors}

\textbf{Definition 2.1.} The ring of truncated Witt vectors
\(W_n(\mathbb{F}_p)\) consists of \(n\)-tuples \((a_0, a_1, \ldots, a_{n-1})\)
with \(a_i \in \mathbb{F}_p\), equipped with addition and multiplication defined
by the \emph{ghost map}:

\[w_k(a) = \sum_{i=0}^{k} p^i a_i^{p^{k-i}} \quad (k = 0, \ldots, n-1)\]

The ghost map \(w: W_n(\mathbb{F}_p) \to \mathbb{Z}^n\) is a ring homomorphism
when the target has componentwise operations:

\[w_k(a + b) = w_k(a) + w_k(b), \quad w_k(a \cdot b) = w_k(a) \cdot w_k(b)\]

The Witt addition and multiplication polynomials \(S_k\), \(P_k\) are obtained
by inverting the ghost map, which requires division by \(p^k\) --- the source of
carries.

\textbf{Definition 2.2} (Frobenius). The Frobenius endomorphism
\(F: W_n(\mathbb{F}_p) \to W_n(\mathbb{F}_p)\) acts by
\(F(a_0, \ldots, a_{n-1}) = (a_0^p, \ldots, a_{n-1}^p)\). Over \(\mathbb{F}_p\),
Fermat's little theorem gives \(a^p = a\), so \(F = \mathrm{id}\) on components.
However, on ghost components: \(w_k(F(a)) = w_{k+1}(a)\) (a shift), which is the
source of the non-trivial arithmetic.

\subsubsection{2.2 The Witt Carry}\label{the-witt-carry}

\textbf{Definition 2.3} (Level-\(k\) Witt carry). For
\(x, y \in W_n(\mathbb{F}_p)\) and an operation
\(\star \in \lbrace{}+, \times\rbrace{}\), the level-\(k\) Witt carry is

\[c_k(x, y) \equiv \frac{w_k(x) \star w_k(y) - \sum_{i=0}^{k-1} p^i z_i^{p^{k-i}}}{p^k} \pmod{p}\]

where \(z = x \star y\) in \(W_n(\mathbb{F}_p)\). This is well-defined because
the ghost map inversion guarantees that the numerator is divisible by \(p^k\).

For \(\star = \times\) (multiplication), the level-2 carry with
\(x = (x_0, x_1, 0)\), \(y = (y_0, y_1, 0) \in W_3(\mathbb{F}_p)\) is:

\[c_2(x, y) \equiv \frac{(x_0^{p^2} + p x_1^p)(y_0^{p^2} + p y_1^p) - z_0^{p^2} - p z_1^p}{p^2} \pmod{p}\]

\subsubsection{2.3 The Fermat Quotient}\label{the-fermat-quotient}

The leading term in the level-2 multiplicative carry is the \emph{Fermat
quotient of the product}:

\[Q_p(x_0, y_0) = \frac{x_0^{p^2} y_0^{p^2} - (x_0 y_0 \bmod p)^{p^2}}{p^2}\]

This is not a random polynomial. Writing \(a = x_0 y_0 \bmod p\) and
\(x_0^p = x_0\), \(y_0^p = y_0\) in \(\mathbb{F}_p\), the quotient reduces to
the classical Fermat quotient

\[q_p(x, y) = \frac{(xy)^p - (xy \bmod p)^p}{p}\]

which is intimately connected to the arithmetic of the Fermat curve
\(X^p + Y^p = Z^p\) and to the Jacobi sums

\[J(\chi_a, \chi_b) = \sum_{x \in \mathbb{F}_p} \chi_a(x) \chi_b(1-x)\]

whose modulus is exactly \(\lvert J\rvert = \sqrt{p}\) --- the Weil bound.

\subsubsection{2.4 The Carry Transfer
Operator}\label{the-carry-transfer-operator}

\textbf{Definition 2.4.} Fix a prime \(p\) and let \(\omega = e^{2\pi i/p}\) be
a primitive \(p\)-th root of unity. The \emph{multiplicative Witt carry matrix}
is the \(p^2 \times p^2\) matrix \(M\) indexed by
\(\mathbb{F}_p \times \mathbb{F}_p\):

\[M_{(x_1,y_1),(x_0,y_0)} = \omega^{c_2(x_0, y_0, x_1, y_1)}\]

where \(c_2\) is computed from the Witt multiplication of
\((x_0, x_1, 0) \times (y_0, y_1, 0)\) in \(W_3(\mathbb{F}_p)\).

The \emph{carry operator} is \(T = M / p^2\).

\textbf{Why the additive character.} The exponentiation
\(c_2 \mapsto \omega^{c_2}\) is not an arbitrary choice but the canonical
discrete Fourier transform on the carry space \(\mathbb{F}_p\). The map
\(\psi: \mathbb{F}_p \to \mathbb{C}^\times\) defined by
\(\psi(t) = e^{2\pi i t/p}\) is the standard additive character of
\(\mathbb{F}_p\), and it is exactly this character that appears in the
construction of Gauss sums
\(\mathfrak{g}(\chi) = \sum_{t \in \mathbb{F}_p} \chi(t) \psi(t)\) and Jacobi
sums \(J(\chi_a, \chi_b) = \sum_x \chi_a(x) \chi_b(1-x)\). These sums compute
the Frobenius eigenvalues on Fermat curves {[}Wei49, IR90{]}.

By weighting each carry transition by \(\psi(c_2) = \omega^{c_2}\), the matrix
\(M\) performs a Fourier analysis of the carry distribution: the eigenvalues of
\(M\) are the spectral components of the carry propagation, decomposed along the
characters of \(\mathbb{F}_p\). The Frobenius eigenvalues emerge because this
Fourier decomposition is the same operation that extracts Jacobi sums from the
point-counting of Fermat varieties.

\textbf{Remark.} The matrix \(M\) also encodes the Frobenius shift: the ``from''
state \((x_0, y_0)\) represents level-0 data, the ``to'' state \((x_1, y_1)\)
represents level-1 data (which becomes level-0 under the Frobenius shift on
ghost components).

\begin{center}\rule{0.5\linewidth}{0.5pt}\end{center}

\newpage

\subsection{3. Proof of Theorem A}\label{proof-of-theorem-a}

\subsubsection{3.1 The Exponent Matrix}\label{the-exponent-matrix}

For \(p = 3\), we enumerate all \(81\) quadruples
\((x_0, y_0, x_1, y_1) \in \mathbb{F}_3^4\), compute the Witt multiplication
\((x_0, x_1, 0) \times_W (y_0, y_1, 0)\) in \(W_3(\mathbb{F}_3)\), and extract
\(c_2 \bmod 3\).

The \(9 \times 9\) exponent matrix \(E\) (where \(M_{j,i} = \omega^{E_{j,i}}\))
has: - 49 entries equal to 0 (carry \(\equiv 0\)) - 16 entries equal to 1 (carry
\(\equiv 1\)) - 16 entries equal to 2 (carry \(\equiv 2\))

The non-trivial entries (carry \(\neq 0\)) are concentrated in rows and columns
corresponding to states \((x, y)\) with \(x, y \in \lbrace{}1, 2\rbrace{}\) (the
units of \(\mathbb{F}_3\)). When either \(x_0 = 0\) or \(y_0 = 0\), the Witt
product has \(z_0 = 0\) and the level-2 carry vanishes.

\subsubsection{3.2 Rationality of the Characteristic
Polynomial}\label{rationality-of-the-characteristic-polynomial}

\textbf{Proposition 3.2.} All coefficients of \(\chi_M(\mu) = \det(\mu I - M)\)
are rational integers.

\emph{Proof.} We compute \(\chi_M\) via the Faddeev--LeVerrier recurrence, which
expresses the coefficients \(c_{n-k}\) of the characteristic polynomial in terms
of the traces \(\mathrm{Tr}(M^k)\). Working over \(\mathbb{Z}[\omega]\) with the
relation \(\omega^2 + \omega + 1 = 0\), we represent each element as
\(a + b\omega\) with \(a, b \in \mathbb{Q}\).

Direct computation shows \(\mathrm{Tr}(M^k) \in \mathbb{Z}\) for all
\(k = 1, \ldots, 9\):

\begin{longtable}[]{@{}ll@{}}
\toprule\noalign{}
\(k\) & \(\mathrm{Tr}(M^k)\) \\
\midrule\noalign{}
\endhead
\bottomrule\noalign{}
\endlastfoot
1 & 3 \\
2 & 15 \\
3 & 135 \\
4 & 927 \\
5 & 4563 \\
6 & 22005 \\
7 & 120123 \\
8 & 659583 \\
9 & 3510135 \\
\end{longtable}

Since each \(\mathrm{Tr}(M^k)\) is rational, the Faddeev--LeVerrier recurrence
produces rational coefficients \(c_j\) at every step. Moreover, the entries of
\(M\) lie in \(\mathbb{Z}[\omega]\), so \(M\) is an algebraic integer matrix and
its eigenvalues are algebraic integers; consequently the coefficients \(c_j\)
(which are elementary symmetric polynomials in the eigenvalues) are algebraic
integers. Being both rational and algebraic integers, they lie in
\(\mathbb{Z}\). \(\square\)

\textbf{Remark} (Galois symmetry). The rationality of the traces is not obvious
\emph{a priori}: the matrix \(M\) has entries in
\(\mathbb{Z}[\omega] \setminus \mathbb{Z}\), and yet
\(\mathrm{Tr}(M^k) \in \mathbb{Z}\) for every \(k\). This is not a numerical
coincidence but reflects a genuine

\[\mathrm{Gal}(\mathbb{Q}(\omega)/\mathbb{Q})\]

-symmetry of the carry operator.

The Galois automorphism \(\sigma: \omega \mapsto \omega^2\) acts on the matrix
\(M\) by replacing \(\omega^{c_2}\) with \(\omega^{2c_2}\). Since the Witt carry
values \(c_2\) take each residue class \(\lbrace{}0, 1, 2\rbrace{}\) with
multiplicities \(49, 16, 16\) (and \(16 = 16\): the non-zero classes appear
equally often), the carry distribution is invariant under the permutation
\(c \mapsto 2c \pmod{3}\) induced by \(\sigma\). This forces \(\sigma(M)\) to be
conjugate to \(M\) by a permutation matrix, so

\[\mathrm{Tr}(\sigma(M)^k) = \mathrm{Tr}(M^k)\]

. Since the only element of \(\mathbb{Q}(\omega)\) fixed by \(\sigma\) is
\(\mathbb{Q}\), the traces are rational.

This Galois invariance is the hallmark of an operator that ``counts points on a
variety over \(\mathbb{F}_p\)'': the Lefschetz trace formula guarantees

\[\sum_i (-1)^i \mathrm{Tr}(\mathrm{Frob}^k \mid H^i) \in \mathbb{Z}\]

for any smooth projective variety, because the point counts
\(N_V(\mathbb{F}_{p^k})\) are manifestly integers. The rationality of
\(\mathrm{Tr}(M^k)\) is thus the first structural signal that \(M\) encodes the
cohomology of an algebraic variety.

\subsubsection{3.3 Factorization}\label{factorization}

The characteristic polynomial is:

\[\chi_M(\mu) = \mu^9 - 3\mu^8 - 3\mu^7 - 27\mu^6 - 99\mu^5 - 135\mu^4\]

Factoring over \(\mathbb{Z}\):

\[\chi_M(\mu) = \mu^4 \cdot (\mu^2 + 3\mu + 3) \cdot (\mu^3 - 6\mu^2 + 12\mu - 45)\]

This factorization is verified by polynomial division: \(\mu^2 + 3\mu + 3\)
divides \(\chi_M(\mu)/\mu^4\) with zero remainder.

\subsubsection{3.4 Identification with the Frobenius
Polynomial}\label{identification-with-the-frobenius-polynomial}

The quadratic \(\mu^2 + 3\mu + 3\) has roots

\[\mu = \frac{-3 \pm i\sqrt{3}}{2}\]

with \(\lvert\mu\rvert^2 = 9/4 + 3/4 = 3\). Hence
\(\lvert\mu\rvert = \sqrt{3} = \sqrt{p}\), the Weil bound for weight 1.

This is the characteristic polynomial of Frobenius \(t^2 - a_p t + p\) with
\(a_p = -3\) (trace of Frobenius). The elliptic curve with this Frobenius at
\(p = 3\) is determined up to isogeny by \(a_3 = -3\).

\textbf{Proposition 3.4.} The curve \(E: y^2 = x^3 + 2x + 1\) over
\(\mathbb{F}_3\) has \(N_E(\mathbb{F}_3) = 7\) and Frobenius characteristic
polynomial \(t^2 + 3t + 3\).

\emph{Proof.} Direct enumeration:

\[E(\mathbb{F}_3) = \lbrace{}(0, \pm 1), (1, \pm 1), (2, \pm 1), \mathcal{O}\rbrace{},\]

giving \(7\) points. Hence \(a_3 = 3 + 1 - 7 = -3\), and the Frobenius
polynomial is \(t^2 - (-3)t + 3 = t^2 + 3t + 3\). Since
\(a_3 \equiv 0 \pmod{3}\), the curve is supersingular. \(\square\)

\subsubsection{3.5 The Additive Case (Negative
Result)}\label{the-additive-case-negative-result}

\textbf{Theorem 3.5} (computer-assisted, exact arithmetic). The additive Witt
carry matrix at \(p = 3\) has characteristic polynomial

\[\chi_{M_{\text{add}}}(\mu) = \mu^6 \cdot (\mu^3 - 6\mu^2 + 9)\]

\emph{which contains no elliptic curve Frobenius factor.}

\emph{Proof.} Same method as Theorem A, replacing Witt multiplication by Witt
addition. The cubic \(\mu^3 - 6\mu^2 + 9\) has three real roots
\(\lbrace{}5.725, 1.399, -1.124\rbrace{}\) (none on
\(\lvert\mu\rvert = \sqrt{3}\)). \(\square\)

This confirms that the Frobenius emerges specifically from the multiplicative
structure, consistent with the fact that the Frobenius endomorphism on Witt
vectors interacts non-trivially only with multiplication.

\begin{center}\rule{0.5\linewidth}{0.5pt}\end{center}

\newpage

\subsection{4. The Lefschetz Trace Formula}\label{the-lefschetz-trace-formula}

\subsubsection{4.1 Point Counting}\label{point-counting}

The Frobenius eigenvalues \(\alpha = (-3 + i\sqrt{3})/2\),
\(\bar{\alpha} = (-3 - i\sqrt{3})/2\) determine the point counts of \(E\) over
all extensions:

\[N_E(\mathbb{F}_{3^k}) = 3^k + 1 - (\alpha^k + \bar{\alpha}^k)\]

\begin{longtable}[]{@{}lll@{}}
\toprule\noalign{}
\(k\) & \(\alpha^k + \bar{\alpha}^k\) & \(N_E(\mathbb{F}_{3^k})\) \\
\midrule\noalign{}
\endhead
\bottomrule\noalign{}
\endlastfoot
1 & \(-3\) & 7 \\
2 & \(+3\) & 7 \\
3 & \(0\) & 28 \\
4 & \(-9\) & 91 \\
5 & \(+27\) & 217 \\
6 & \(-54\) & 784 \\
\end{longtable}

\subsubsection{4.2 The Local Zeta Function}\label{the-local-zeta-function}

The local zeta function of \(E/\mathbb{F}_3\) is:

\[Z(E/\mathbb{F}_3, T) = \frac{1 + 3T + 3T^2}{(1 - T)(1 - 3T)}\]

The numerator \(1 + 3T + 3T^2 = (1 - \alpha T)(1 - \bar{\alpha} T)\) with
\(\lvert\alpha\rvert = \sqrt{3}\), confirming the Riemann Hypothesis for this
curve.

\subsubsection{4.3 Connection to the Carry Operator
Trace}\label{connection-to-the-carry-operator-trace}

The traces of \(M^k\) encode a richer object than the Frobenius pair alone (they
include contributions from the cubic factor and the zero eigenvalues). The
Frobenius contribution to \(\mathrm{Tr}(M^k)\) is \(\alpha^k + \bar{\alpha}^k\),
which can be extracted via the Weil zeta function decomposition.

\begin{center}\rule{0.5\linewidth}{0.5pt}\end{center}

\newpage

\subsection{5. Computational Evidence and
Extensions}\label{computational-evidence-and-extensions}

\subsubsection{\texorpdfstring{5.1 The Additive Case at \(p = 2\): The Golden
Ratio and the Ihara
Zeta}{5.1 The Additive Case at p = 2: The Golden Ratio and the Ihara Zeta}}\label{the-additive-case-at-p-2-the-golden-ratio-and-the-ihara-zeta}

For \(p = 2\), the additive Witt carry operator on \(W_3(\mathbb{F}_2)\) has two
non-zero eigenvalues of \(T = M/4\):

\[\lambda = \frac{1 \pm \sqrt{5}}{4}\]

satisfying \(4\lambda^2 - 2\lambda - 1 = 0\). The dominant eigenvalue is
\(\lambda_0 = \varphi/2\), where \(\varphi = (1+\sqrt{5})/2\) is the golden
ratio.

This is not a curiosity. In base 2, the Fermat curve \(X^2 + Y^2 = 1\) is
trivial (genus 0), and the ``elliptic curve'' degenerates. What remains is a
\emph{combinatorial} structure: the carry propagation in binary addition is
equivalent to the random walk on the Fibonacci graph
\(\bullet \rightleftharpoons \bullet\), whose adjacency matrix has eigenvalues
\(\pm \varphi\).

The connection to spectral graph theory is precise. The \emph{Ihara zeta
function} of a \((q+1)\)-regular graph \(G\) is

\[Z_G(u) = \prod_{\gamma} (1 - u^{|\gamma|})^{-1}\]

where the product runs over primitive cycles. For the simplest non-trivial graph
(the loop of length 2), the Ihara zeta function has poles governed by
\(x^2 - x - 1 = 0\) --- the minimal polynomial of \(\varphi\). The carry
operator at \(p = 2\) is thus computing the Ihara zeta of the carry propagation
graph, rather than the Hasse--Weil zeta of a variety.

This degeneration is consistent with the genus formula: \(g(C_2) = 0\), so there
are no Frobenius eigenvalues in the Weil sense. The golden ratio is the
graph-theoretic shadow of the absent algebraic geometry. Notably, Ramanujan
graphs --- the optimal expanders --- are defined by the condition that their
non-trivial Ihara eigenvalues satisfy \(\lvert\lambda\rvert \leq 2\sqrt{q}\),
the graph-theoretic analogue of the Weil bound. The carry operator at \(p = 2\)
saturates this bound, suggesting that carry propagation graphs are Ramanujan.

\subsubsection{5.2 Cross-Prime Evidence}\label{cross-prime-evidence}

For primes \(p = 2, 3, 5, 7, 11, 13\), the multiplicative Witt carry operator at
level 2 with state \((x_0, y_0)\) was computed:

\begin{longtable}[]{@{}
  >{\raggedright\arraybackslash}p{(\columnwidth - 6\tabcolsep) * \real{0.0943}}
  >{\raggedright\arraybackslash}p{(\columnwidth - 6\tabcolsep) * \real{0.1321}}
  >{\raggedright\arraybackslash}p{(\columnwidth - 6\tabcolsep) * \real{0.6415}}
  >{\raggedright\arraybackslash}p{(\columnwidth - 6\tabcolsep) * \real{0.1321}}@{}}
\toprule\noalign{}
\begin{minipage}[b]{\linewidth}\raggedright
\(p\)
\end{minipage} & \begin{minipage}[b]{\linewidth}\raggedright
\(p^2\)
\end{minipage} & \begin{minipage}[b]{\linewidth}\raggedright
Frobenius at \(\lvert\mu\rvert = \sqrt{p}\)?
\end{minipage} & \begin{minipage}[b]{\linewidth}\raggedright
Notes
\end{minipage} \\
\midrule\noalign{}
\endhead
\bottomrule\noalign{}
\endlastfoot
2 & 4 & No (\(\lvert\mu\rvert \neq \sqrt{2}\)) & Genus 0; Ihara zeta instead \\
3 & 9 & \textbf{Yes --- exact} (Theorem A) & \(\mu^2 + 3\mu + 3\) \\
5 & 25 & No (untwisted) & Expected: degree 12 (genus 6); see §5.4 \\
7 & 49 & \textbf{Yes --- exact} (6 eigenvalues) & See §5.3 \\
11 & 121 & No & Expected: degree 90 (genus 45) \\
13 & 169 & No & Expected: degree 132 (genus 66) \\
\end{longtable}

The exactness at \(p = 3\) is attributed to the fact that the Fermat quotient
has only one non-trivial case (\((x_0, y_0) = (2,2)\)), producing a clean
algebraic structure. For most \(p \geq 5\), multiple non-trivial cases create
interference. The notable exception is \(p = 7\), where 6 eigenvalues at
\(\lvert\mu\rvert = \sqrt{7}\) emerge despite the Fermat curve having genus 15.

\subsubsection{\texorpdfstring{5.3 The Second Frobenius: Exact Factorization at
\(p = 7\) (Theorem
B)}{5.3 The Second Frobenius: Exact Factorization at p = 7 (Theorem B)}}\label{the-second-frobenius-exact-factorization-at-p-7-theorem-b}

The \(49 \times 49\) Witt carry matrix for \(p = 7\) was computed exactly over
\(\mathbb{Z}[\omega_7]\) using the Faddeev--LeVerrier algorithm. All
characteristic polynomial coefficients are rational integers, confirming the
Galois symmetry observed at \(p = 3\).

\textbf{Theorem B.} The Weil-bound eigenvalues of the \(49 \times 49\) carry
matrix at \(p = 7\) form a genus-3 Frobenius factor:

\[\chi_M(\mu) = \mu^{14} \cdot \underbrace{(\mu^2 - 7)^2 \cdot (\mu^2 + 7)}_{\text{Weil factor}} \cdot (\mu^5 - 35\mu^3 - 49\mu^2 + 147\mu - 343)^2 \cdot (\mu^6 - \cdots)^2 \cdot (\mu^7 - \cdots)\]

(The full degree-49 factorization is available in the computational scripts
\texttt{F06\_witt\_algebra.py} and \texttt{F08c\_p7\_factorization.py} in the
experiments directory; the remaining factors have no eigenvalues at the Weil
bound.)

The Weil-bound factor is:

\[(\mu^2 - 7)^2 \cdot (\mu^2 + 7) = \mu^6 - 7\mu^4 - 49\mu^2 + 343\]

This decomposes as:

\begin{itemize}
\tightlist
\item
  \((\mu^2 + 7)\): the \textbf{Frobenius polynomial of a supersingular elliptic
  curve} with \(a_7 = 0\), satisfying \(N_E(\mathbb{F}_7) = 8\).
\item
  \((\mu^2 - 7)^2\): a \textbf{genus-2 superspecial Frobenius polynomial} (roots
  \(\pm\sqrt{7}\), each with multiplicity 2; product \(= 49 = 7^2\)).
\end{itemize}

The combined degree-6 factor satisfies the \textbf{functional equation} for
genus-3 Frobenius: \(7^3 P(\mu) = \mu^6 P(7/\mu)\). This confirms that 6 of the
30 expected Frobenius eigenvalues of the genus-15 Fermat curve \(C_7\) are
captured by the \(49 \times 49\) carry matrix.

\textbf{Remark.} The 6 eigenvalues at \(\lvert\mu\rvert = \sqrt{7}\) are a
\emph{subset} of the full Frobenius spectrum of \(C_7\). The remaining 24
eigenvalues presumably require a deeper Witt truncation (\(n \geq 3\)) to
resolve, consistent with the genus obstruction (§6.1).

\subsubsection{\texorpdfstring{5.4 Dirichlet Character Twist: Rescuing
\(p = 5\)}{5.4 Dirichlet Character Twist: Rescuing p = 5}}\label{dirichlet-character-twist-rescuing-p-5}

For primes where the untwisted carry matrix yields no Weil-bound eigenvalues, a
natural question is whether twisting by a Dirichlet character can extract the
missing Frobenius.

Define the twisted matrix:

\[M(\chi)_{(x_1,y_1),(x_0,y_0)} = \chi(x_0 y_0 \bmod p) \cdot \omega^{c_2}\]

Three twist variants were tested: multiplicative (\(\chi(x_0 y_0)\)), additive
(\(\chi(x_0 + y_0)\)), and carry-value (\(\chi(c_2) \cdot \omega^{c_2}\)).

Result at \(p = 5\): The untwisted matrix has no eigenvalues at
\(\lvert\mu\rvert = \sqrt{5}\). However, the additive twist with \(\chi_2\) (the
Legendre symbol mod 5, order 2) produces 2 eigenvalues at exactly
\(\lvert\mu\rvert = \sqrt{5} = 2.236068\):

\[\mu = +2.0237 - 0.9511i, \quad \mu = -2.0237 - 0.9511i\]

Their product is \(\mu_1 \mu_2 = -5\) (anti-conjugate pair, not a standard
Frobenius conjugate pair with product \(+p\)).

Result at \(p = 7\): All 6 multiplicative twists (\(\chi_0\) through \(\chi_5\))
produce exactly 6 eigenvalues at \(\lvert\mu\rvert = \sqrt{7}\), each with
distinct phases. The trivial and \(\chi_3\) (order 2) twists give degenerate
phases (multiples of \(\pi/2\)); the remaining twists produce non-trivial phases
consistent with Jacobi sum arguments.

This decomposition by character is consistent with the Jacobi sum structure of
the Frobenius on Fermat curves, where each character pair \((\chi_a, \chi_b)\)
selects a specific eigenvalue.

\subsubsection{5.5 Level-3 and Level-4 Marginals (Negative
Result)}\label{level-3-and-level-4-marginals-negative-result}

An attempt was made to deepen the carry level while keeping the
\(p^2 \times p^2\) state space, by marginalizing over auxiliary Witt components:

\[M_{(x_1,y_1),(x_0,y_0)} = \sum_{x_2, y_2} \omega^{c_3(x_0,x_1,x_2,0;\, y_0,y_1,y_2,0)}\]

\textbf{Result.} At all primes tested (\(p = 3, 5, 7\)), the level-3 and level-4
marginal matrices collapse to \textbf{rank 1}: a single nonzero eigenvalue at
\(p^{k+1}\) (e.g., \(625 = 5^4\) for \(p = 5\), level 3), with all other
eigenvalues zero.

The root cause is \textbf{character orthogonality}: the sum
\(\sum_{x_2,y_2} \omega^{c_k}\) evaluates to \(p^2\) when a constraint is met
and \(0\) otherwise, destroying all spectral structure.

\textbf{Conclusion.} Level 2 is the \textbf{only productive level} for the dense
marginal approach, because it involves no marginalization --- each matrix entry
is a single \(\omega^{c_2}\).

\subsubsection{\texorpdfstring{5.6 Gauss Sum Identification at
\(p = 7\)}{5.6 Gauss Sum Identification at p = 7}}\label{gauss-sum-identification-at-p-7}

The Frobenius eigenvalues of the carry matrix were compared against all Gauss
sums \(g(\chi^k) = \sum_{x \in \mathbb{F}_p^{\ast}} \chi^k(x) \omega^x\) and
Jacobi sums \(J(\chi^a, \chi^b) = \sum_x \chi^a(x) \chi^b(1-x)\) for
multiplicative characters on \(\mathbb{F}_p^{\ast}\).

\textbf{Result.} At \(p = 7\), the supersingular eigenvalues
\(\mu = \pm i\sqrt{7}\) from the factor \((\mu^2 + 7)\) match the quadratic
Gauss sum \textbf{exactly}:

\[\mu = +i\sqrt{7} = g(\chi^3), \quad \mu = -i\sqrt{7} = -g(\chi^3)\]

where \(\chi^3\) is the Legendre symbol (quadratic character) mod 7, and
\(g(\chi^3) = \sum_{x=1}^{6} \left(\frac{x}{7}\right) \omega^x = i\sqrt{7}\).

The superspecial eigenvalues \(\mu = \pm\sqrt{7}\) from \((\mu^2 - 7)^2\) do
\textbf{not} match any standard Gauss or Jacobi sum. At \(p = 3\), the Frobenius
eigenvalues \((-3 \pm i\sqrt{3})/2\) (with phases \(\pm 5\pi/6\)) likewise do
not match the sole Gauss sum \(g(\chi) = i\sqrt{3}\) (phase \(\pi/2\)).

\textbf{Interpretation.} The Fermat curve \(X^p + Y^p = 1\) degenerates in
characteristic \(p\) (becoming \(X + Y = 1\), a line), so the standard Jacobi
sum formula for Fermat curve Frobenius does not apply. The carry eigenvalues
arise from a different variety --- the Artin-Schreier variety
\(z^p - z = c_2(x_0, x_1, y_0, y_1)\) --- whose Frobenius structure only
partially overlaps with the classical Gauss/Jacobi catalog.

\subsubsection{5.7 Structural Properties of the Carry
Exponent}\label{structural-properties-of-the-carry-exponent}

The distribution of \(c_2\) values over \(\mathbb{F}_p^4\) reveals a striking
balance: nonzero values are \textbf{exactly equidistributed}.

\begin{longtable}[]{@{}
  >{\raggedright\arraybackslash}p{(\columnwidth - 6\tabcolsep) * \real{0.0847}}
  >{\raggedleft\arraybackslash}p{(\columnwidth - 6\tabcolsep) * \real{0.3390}}
  >{\raggedleft\arraybackslash}p{(\columnwidth - 6\tabcolsep) * \real{0.4576}}
  >{\raggedright\arraybackslash}p{(\columnwidth - 6\tabcolsep) * \real{0.1186}}@{}}
\toprule\noalign{}
\begin{minipage}[b]{\linewidth}\raggedright
\(p\)
\end{minipage} & \begin{minipage}[b]{\linewidth}\raggedleft
\(N_0\) (\(c_2 = 0\))
\end{minipage} & \begin{minipage}[b]{\linewidth}\raggedleft
\(N_a\) (\(c_2 = a \neq 0\))
\end{minipage} & \begin{minipage}[b]{\linewidth}\raggedright
Check
\end{minipage} \\
\midrule\noalign{}
\endhead
\bottomrule\noalign{}
\endlastfoot
3 & 49 & 16 each & \(49 + 2 \cdot 16 = 81 = 3^4\) \\
5 & 273 & 88 each & \(273 + 4 \cdot 88 = 625 = 5^4\) \\
7 & 745 & 276 each & \(745 + 6 \cdot 276 = 2401 = 7^4\) \\
\end{longtable}

Consequently, all exponential sums
\(S_a = \sum_{\boldsymbol{x} \in \mathbb{F}_p^4} \omega^{a \cdot c_2(\boldsymbol{x})}\)
for \(a \neq 0\) are \textbf{real and equal}: \(S_a = (p N_0 - p^4) / (p - 1)\).

A deeper structural property: the Weil-bound eigenvalues
(\(\lvert\mu\rvert = \sqrt{p}\)) have \textbf{coefficient zero} in the expansion
\(\langle \boldsymbol{1} \mid M^k \mid \boldsymbol{1} \rangle = \sum_i c_i \lambda_i^k\).
The all-ones vector is orthogonal to every Weil eigenvector, so the Frobenius is
``invisible'' in the total exponential sum but fully present in
\(\mathrm{Tr}(M^k)\) (the periodic orbit sum).

\subsubsection{\texorpdfstring{5.8 Extended Twist Scan: \(p = 11\) and
\(p = 13\)}{5.8 Extended Twist Scan: p = 11 and p = 13}}\label{extended-twist-scan-p-11-and-p-13}

The character twist methodology of §5.4 was extended to \(p = 11\) (121 × 121
matrix) and \(p = 13\) (169 × 169 matrix), testing all \(p-1\) multiplicative
and \(p-1\) additive twists.

\(p = 11\): Productive with higher-order characters. A comprehensive scan of all
character twists at \(p = 11\) reveals that additive twists \(\chi_2\) and
\(\chi_8\) (both of order 5) each produce 1 eigenvalue at
\(\lvert\mu\rvert \approx \sqrt{11}\), with traces \(a_{11}(\chi_2) = -0.76\)
and \(a_{11}(\chi_8) = 2.10\). Lower-order characters are not productive at this
prime.

\(p = 13\): Partial. Additive twists \(\chi_4\) and \(\chi_8\) (both of order 3)
each produce 1 eigenvalue at \(\lvert\mu\rvert = \sqrt{13}\):

\[\mu = -1.264 \pm 3.378 i, \quad \lvert\mu\rvert = 3.606 \approx \sqrt{13}\]

These are complex conjugates (related by
\(\chi_4 \leftrightarrow \chi_8 = \overline{\chi_4}\)). No other twist type is
productive.

\textbf{Extended prime scan} (primes up to 23):

\begin{longtable}[]{@{}lrcc@{}}
\toprule\noalign{}
\(p\) & Untwisted & Best productive twist & Total Weil eigenvalues \\
\midrule\noalign{}
\endhead
\bottomrule\noalign{}
\endlastfoot
3 & 2 & mul Legendre: 2 & 4 \\
5 & 0 & add Legendre: 2 & 8 \\
7 & 6 & all mul + add: 6 each & 40+ \\
11 & 0 & add ord 5: 1 each & 2 \\
13 & 0 & add ord 3: 1 each & 2 \\
17 & 0 & add ord 16: 2; ord 8: 3 & 14 \\
19 & 6 & add ord 6,9: 2 each & 23+ \\
23 & 0 & add ord 11: 1--2 each & 14+ \\
\end{longtable}

\textbf{Every prime tested produces Weil eigenvalues with appropriate character
twists.} The productive character order appears to depend on \(p\) in a
non-trivial way (e.g., order 5 at \(p = 11\), order 3 at \(p = 13\), order 16 at
\(p = 17\)). No simple arithmetic formula predicts the productive order. The
anomalous richness at \(p = 7\) and \(p = 19\) (large untwisted Weil counts) and
the requirement for high-order characters at \(p = 11, 17, 23\) remain open
questions.

\begin{center}\rule{0.5\linewidth}{0.5pt}\end{center}

\newpage

\subsection{6. The Artin-Schreier / Fermat Curve
Conjecture}\label{the-artin-schreier-fermat-curve-conjecture}

This is the central conjecture of the paper, motivated by the structure of the
Witt carry, the Gauss sum identification at \(p = 7\), and the failure of naïve
generalization to \(p \geq 5\).

\subsubsection{6.1 The Genus Obstruction}\label{the-genus-obstruction}

For \(p = 3\), the quadratic factor \(\mu^2 + 3\mu + 3\) in \(\chi_M\) naturally
suggests searching for a quadratic Frobenius factor at all primes. However, the
Frobenius polynomial at general \(p\) is not quadratic: the degree grows with
the genus of the Fermat curve.

The resolution lies in the Fermat quotient. The level-2 multiplicative carry is
dominated by

\[Q_p(x, y) = \frac{x^p y^p - (xy \bmod p)^p}{p}\]

This term governs the arithmetic of the \emph{Fermat curve}
\(C_p: X^p + Y^p = 1\) over \(\mathbb{F}_p\). By the genus formula for smooth
plane curves:

\[g(C_p) = \frac{(p-1)(p-2)}{2}\]

The Weil conjectures (proved by Deligne {[}Del74{]}) guarantee that the zeta
function numerator of \(C_p\) is a polynomial of degree \(2g = (p-1)(p-2)\) in
\(T\), with all roots on \(\lvert T\rvert = p^{-1/2}\).

\begin{longtable}[]{@{}llll@{}}
\toprule\noalign{}
\(p\) & \(g(C_p)\) & \(2g\) & Frobenius polynomial degree \\
\midrule\noalign{}
\endhead
\bottomrule\noalign{}
\endlastfoot
2 & 0 & 0 & (trivial) \\
3 & 1 & 2 & \textbf{Quadratic} (Theorem A) \\
5 & 6 & 12 & Degree 12 \\
7 & 15 & 30 & Degree 30 \\
11 & 45 & 90 & Degree 90 \\
\end{longtable}

For \(p = 3\), genus 1 means \(C_3\) is an elliptic curve, and the Frobenius
polynomial is quadratic --- which is precisely the factor we proved. For
\(p = 5\), the expected Frobenius polynomial has degree 12. One cannot find it
in the characteristic polynomial of a \(25 \times 25\) matrix (degree 25), which
barely exceeds \(2g = 12\), and certainly not as a quadratic factor.

\subsubsection{6.2 Jacobi Sum Connection}\label{jacobi-sum-connection}

The Frobenius eigenvalues on the Fermat curve \(C_p\) are expressed by Jacobi
sums. For characters \(\chi_a: \mathbb{F}_p^\times \to \mathbb{C}^\times\) of
order \(p-1\):

\[J(\chi_a, \chi_b) = \sum_{x \in \mathbb{F}_p} \chi_a(x) \chi_b(1-x)\]

The \(2g\) Frobenius eigenvalues of \(C_p\) are exactly the Jacobi sums
\(J(\chi_a, \chi_b)\) for appropriate pairs \((a, b)\), each satisfying
\(\lvert J\rvert = \sqrt{p}\).

The Witt carry matrix \(M\) ``sees'' these Jacobi sums because the additive
character \(\omega^{c_2}\) in the matrix entries performs a discrete Fourier
transform over the carry values, and the carry values themselves are determined
by the Fermat quotient --- which is the polynomial identity underlying the
Jacobi sum.

\textbf{Conjecture 1} (Artin-Schreier Carry Conjecture). \emph{The Witt carry
operator at level 2 computes (a subfactor of) the Frobenius on the
Artin-Schreier variety}

\[V_p:\ z^p - z = c_2(x_0, x_1, y_0, y_1) \quad \text{over } \mathbb{F}_p,\]

where \(c_2\) is the level-2 multiplicative Witt carry. This variety is
dominated by the Fermat quotient

\[\frac{x^p y^p - (xy)^p}{p}.\]

\emph{Specifically:}

\begin{enumerate}
\def\labelenumi{(\alph{enumi})}
\tightlist
\item
  For every prime \(p\), the characteristic polynomial of the multiplicative
  Witt carry matrix \(M\) contains, as an exact factor, a polynomial
  \(P_p(\mu)\) whose roots lie on
\end{enumerate}

\[\lvert \mu \rvert = \sqrt{p},\]

and \(P_p(\mu)\) divides the Frobenius polynomial of \(V_p\).

\begin{enumerate}
\def\labelenumi{(\alph{enumi})}
\setcounter{enumi}{1}
\item
  For \(p = 3\), \(P_3(\mu) = \mu^2 + 3\mu + 3\) is the full Frobenius
  polynomial of \(C_3\) (genus 1, Theorem A).
\item
  For \(p = 7\),
\end{enumerate}

\[P_7(\mu) = (\mu^2 - 7)^2(\mu^2 + 7)\]

captures 6 Frobenius eigenvalues (Theorem B). The factor \(\mu^2 + 7\) has roots
equal to the quadratic Gauss sum

\[g(\chi^3) = i\sqrt{7}\]

\emph{(§5.6).}

\begin{enumerate}
\def\labelenumi{(\alph{enumi})}
\setcounter{enumi}{3}
\item
  For \(p \geq 5\) in general, the Frobenius polynomial of the underlying
  variety has degree \(\leq 2g = (p-1)(p-2)\). The \(p^2 \times p^2\) carry
  matrix captures a subfactor of bounded degree. The complete Frobenius requires
  either a deeper Witt truncation (\(n \geq 3\)), or a Dirichlet character
  decomposition (§5.4), or both.
\item
  The carry exponent \(c_2\) is equidistributed on \(\mathbb{F}_p^{\ast}\)
  (§5.7) and the Weil eigenvectors are orthogonal to \(\lvert 1 \rangle\),
  consistent with the Artin-Schreier interpretation where the Frobenius lives in
  the non-trivial additive character sector.
\end{enumerate}

\subsubsection{\texorpdfstring{6.3 The Case \(p = 5\): Character
Decomposition}{6.3 The Case p = 5: Character Decomposition}}\label{the-case-p-5-character-decomposition}

For \(p = 5\), the untwisted \(25 \times 25\) carry matrix yields no eigenvalues
at \(\lvert\mu\rvert = \sqrt{5}\). This is consistent with the genus
obstruction: the full Frobenius polynomial of \(C_5\) has degree 12, and the
untwisted matrix cannot accommodate it.

However, twisting by the additive Legendre character \(\chi_2\) extracts 2
eigenvalues at exactly \(\lvert\mu\rvert = \sqrt{5}\) (§5.4). This suggests that
the Frobenius is present but ``hidden'' in the character decomposition of the
carry matrix, analogous to how the Frobenius eigenvalues of the Fermat curve
decompose by Jacobi sum index \((\chi_a, \chi_b)\).

The Frobenius eigenvalues of \(C_5: X^5 + Y^5 = 1\) are the 12 Jacobi sums
\(J(\chi_a, \chi_b)\) for characters of order 4 on \(\mathbb{F}_5^\times\).
These are well-tabulated (see {[}IR90, Ch. 8{]}). The full program requires:

\begin{enumerate}
\def\labelenumi{\arabic{enumi}.}
\tightlist
\item
  Computing the \(25 \times 25\) carry matrix twisted by all \((p-1)^2 = 16\)
  pairs of characters.
\item
  Extracting the Weil-bound eigenvalues from each twist.
\item
  Matching the resulting eigenvalues against the tabulated Jacobi sums.
\end{enumerate}

If the union of Weil-bound eigenvalues over all twists reproduces the 12
Frobenius eigenvalues of \(C_5\), this would confirm Conjecture 1(d) at
\(p = 5\).

\begin{center}\rule{0.5\linewidth}{0.5pt}\end{center}

\newpage

\subsection{7. The Overlap Transfer Matrix (Negative
Result)}\label{the-overlap-transfer-matrix-negative-result}

\subsubsection{7.1 Construction}\label{construction}

We also studied the \emph{overlapping} transfer matrix on the full state space
\(W_2(\mathbb{F}_p) \times W_2(\mathbb{F}_p)\), of dimension \(p^4\), where the
``from'' state is the complete Witt pair \((x_0, x_1, y_0, y_1)\) and the ``to''
state is the shifted pair \((x_1, x_2, y_1, y_2)\), with the carry weight
\(\omega^{c_2}\).

\subsubsection{7.2 Result: Uniform Modulus}\label{result-uniform-modulus}

\textbf{Proposition 7.2.} For \(p = 2, 3, 5\), all eigenvalues of the overlap
transfer matrix have modulus exactly \(p\).

This is a topological consequence of the shift structure. The overlap operator
is an \emph{expanding map} on a symbolic space with alphabet size \(p^2\), and
the Ruelle--Perron--Frobenius theorem {[}Rue04{]} forces the spectral radius to
equal the topological degree of the map, which is \(p\) (corresponding to \(p\)
choices of new carry input at each step).

\subsubsection{7.3 Phase Preservation}\label{phase-preservation}

Despite the uniform modulus, the \emph{phases} of the eigenvalues are preserved
from the dense (non-overlap) formulation. For \(p = 3\), the overlap matrix has
eigenvalues at arguments \(\pm 5\pi/6\) --- the same Frobenius phase as in
Theorem A. This demonstrates that:

\begin{itemize}
\tightlist
\item
  The \textbf{modulus} of Frobenius eigenvalues is encoded in the \emph{dense}
  carry matrix (Section 3), where the averaging over level-1 digits creates the
  correct contraction.
\item
  The \textbf{phase} of Frobenius eigenvalues is an intrinsic property of the
  additive character weighting \(\omega^{c_2}\), invariant under the choice of
  shift structure.
\end{itemize}

The phase is preserved because it is encoded in the argument of the exponential
\(\exp(2\pi i \cdot c_2/p)\). The constructive/destructive interference of Witt
carries (which defines the Gauss/Jacobi sums) cannot be destroyed by a shift,
because it is an invariant of the character class.

\begin{center}\rule{0.5\linewidth}{0.5pt}\end{center}

\newpage

\subsection{8. Conjectures and Future
Directions}\label{conjectures-and-future-directions}

\subsubsection{8.1 The Universal Carry
Variety}\label{the-universal-carry-variety}

\textbf{Conjecture 2} (Universality via Artin-Schreier/Fermat varieties). The
carry variety \(V_p: z^p - z = c_2(x_0, x_1, y_0, y_1)\) is birationally related
to the Fermat curve \(C_p: X^p + Y^p = 1\) over \(\mathbb{F}_p\). The
characteristic polynomial of Frobenius on \(H^{\bullet}(V_p)\) --- whose
Weil-bound roots have modulus \(\sqrt{p}\) --- divides the characteristic
polynomial of the multiplicative Witt carry operator on the state space
\(W_n(\mathbb{F}_p) \times W_n(\mathbb{F}_p)\) for sufficiently large \(n\).

For \(p = 3\), \(V_3\) reduces to a genus-1 curve whose Frobenius is
\(\mu^2 + 3\mu + 3\) (Theorem A). For \(p = 7\), the Artin-Schreier
interpretation is supported by three structural facts: (i) the Gauss sum
identification \(g(\chi^3) = i\sqrt{7}\) (§5.6), which is the standard Frobenius
eigenvalue of Artin-Schreier covers; (ii) the equidistribution of \(c_2\) on
\(\mathbb{F}_p^{\ast}\) (§5.7), which is the condition for the Artin-Schreier
surface \(z^p - z = c_2\) to have well-defined cohomology in every non-trivial
character sector; and (iii) the orthogonality of Weil eigenvectors to
\(\lvert 1\rangle\) (§5.7), which means the Frobenius lives entirely in the
non-trivial sectors of the Artin-Schreier cover --- the trivial sector being the
``base'' variety \(\mathbb{F}_p^4\).

\subsubsection{8.2 Toward an Euler Product (and a Fundamental
Obstruction)}\label{toward-an-euler-product-and-a-fundamental-obstruction}

If each prime's carry operator produced the local factor of a single global
variety, one could assemble an Euler product \(L(s) = \prod_p Z_p(p^{-s})\).

\textbf{Cross-prime inconsistency.} Computational experiments reveal that the
Witt carry operator at different primes does \textbf{not} encode a single
algebraic curve over \(\mathbb{Q}\). Specifically:

\begin{itemize}
\tightlist
\item
  At \(p = 3\), the Frobenius trace is \(a_3 = -3\) (Theorem A), matching the
  supersingular curve \(E: y^2 = x^3 + 2x + 1\).
\item
  At \(p = 7\), the supersingular Frobenius factor gives \(a_7 = 0\) (§5.3).
\item
  However, \(E/\mathbb{F}_7\) has \(a_7 = 3\), not \(0\).
\end{itemize}

A database search identifies 6 elliptic curves over \(\mathbb{Q}\) satisfying
both \(a_3 = -3\) and \(a_7 = 0\) (e.g., \(y^2 = x^3 - x^2 - x + 1\)), but none
match at all primes tested.

\textbf{Interpretation.} The carry variety \(V_p\) depends on \(p\): the Fermat
curve \(C_p: X^p + Y^p = 1\) is a different curve at each prime, not the
reduction mod \(p\) of a fixed curve over \(\mathbb{Q}\). This is consistent
with the geometric picture --- the carry operator encodes the \emph{local}
arithmetic of the Fermat quotient at each prime, and there is no single global
object of which these are reductions.

The path to a global \(L\)-function, if it exists, must therefore involve
either: - A more sophisticated global construction (e.g., a family
\(\lbrace{}C_p\rbrace{}_p\) parameterized by \(p\), in the spirit of Hasse--Weil
for families of varieties). - A product over the \emph{character-decomposed}
local factors from §5.4, which may exhibit better cross-prime coherence.

\subsubsection{8.3 The Cubic Factor}\label{the-cubic-factor}

The cubic \(\mu^3 - 6\mu^2 + 12\mu - 45\) in \(\chi_M\) has Vieta relations
\(\sigma_1 = 2p\), \(\sigma_2 = p(p+1)\), \(\sigma_3 = p^2(p+2)\), all divisible
by \(p\). The moduli of the complex roots (\(\approx 2.905\)) do not match any
standard Weil bound, suggesting this factor encodes higher-weight or non-smooth
cohomological data. Its identification remains open.

\subsubsection{8.4 Computational Frontiers}\label{computational-frontiers}

The level-3 and level-4 marginal approach (§5.5) has been ruled out:
marginalization destroys spectral structure via character orthogonality, and the
only productive dense level is level 2.

The viable strategies for extending the results are:

\begin{enumerate}
\def\labelenumi{\arabic{enumi}.}
\item
  Character decomposition at \(p = 5\) (§6.3): Compute the \(25 \times 25\)
  carry matrix twisted by all 16 character pairs \((\chi_a, \chi_b)\) and
  collect Weil-bound eigenvalues. This requires exact arithmetic over
  \(\mathbb{Z}[\zeta_5]\).
\item
  Extended state space at \(p = 7\): The current \(49 \times 49\) matrix
  captures 6 of 30 Frobenius eigenvalues. The full state
  \(W_2(\mathbb{F}_7) \times W_2(\mathbb{F}_7)\) (\(2401 \times 2401\)) with
  exact arithmetic over \(\mathbb{Z}[\omega_7]\) should reveal more.
\item
  \textbf{Systematic Jacobi sum matching}: Tabulate the known Jacobi sums
  \(J(\chi_a, \chi_b)\) for \(p = 5, 7\) and compare with the eigenvalues
  obtained from character-twisted carry matrices. This would give a precise
  identification of which Frobenius eigenvalues the carry operator ``sees'' at
  each character.
\end{enumerate}

\begin{center}\rule{0.5\linewidth}{0.5pt}\end{center}

\newpage

\subsection{9. Conclusion}\label{conclusion}

The carry operator of Witt vector multiplication --- an object as elementary as
the ``carrying'' step in schoolbook arithmetic --- encodes Frobenius eigenvalues
of algebraic varieties over finite fields. This is not an analogy: it is an
exact algebraic identity, verified computationally by the Faddeev--LeVerrier
algorithm with exact arithmetic over \(\mathbb{Z}[\omega_p]\).

The principal results are:

\begin{enumerate}
\def\labelenumi{\arabic{enumi}.}
\item
  \textbf{Theorem A} (\(p = 3\)): The characteristic polynomial of the
  \(9 \times 9\) carry matrix contains the Frobenius of the supersingular
  elliptic curve \(E: y^2 = x^3 + 2x + 1\) over \(\mathbb{F}_3\).
\item
  \textbf{Theorem B} (\(p = 7\)): The \(49 \times 49\) carry matrix yields a
  genus-3 Frobenius factor \((\mu^2 - 7)^2(\mu^2 + 7)\) with 6 eigenvalues at
  the Weil bound. The supersingular pair \(\pm i\sqrt{7}\) is exactly the
  quadratic Gauss sum \(g(\chi^3)\) of the Legendre symbol --- the first
  identification of a carry eigenvalue with a named quantity of analytic number
  theory.
\item
  \textbf{Character universality}: Dirichlet character twists of the carry
  operator recover Weil-bound eigenvalues at \textbf{every prime tested}
  (\(p = 3, 5, 7, 11, 13, 17, 19, 23\)). The additive Legendre twist works at
  \(p = 5\) and \(p = 17\); order-3 characters at \(p = 13\); order-5 characters
  at \(p = 11\); higher-order characters at \(p = 17, 23\). The productive
  character order varies with \(p\) in a currently unexplained way.
\item
  \textbf{Structural theorems}: The carry exponent \(c_2\) is equidistributed on
  \(\mathbb{F}_p^{\ast}\), and Weil-bound eigenvectors are orthogonal to
  \(\lvert 1\rangle\) --- explaining why the Frobenius is invisible in the total
  exponential sum but present in the periodic orbit structure
  \(\mathrm{Tr}(M^k)\). This is the spectral signature of an Artin-Schreier
  variety \(z^p - z = c_2(\boldsymbol{x})\).
\item
  \textbf{Binary carry complex spectrum}: The D-odd conditioned binary carry
  chain (from {[}E{]}, whose sector ratio is conjectured to converge to
  \(-\pi\), Conjecture 1 of {[}P1{]}) has \textbf{complex} dominant eigenvalues
  \(\lambda \approx -0.302 \pm 0.340\,i\) (modulus \(\approx 0.454\), phase
  \(\approx 131.6°\)), indicating oscillatory approach to equilibrium. The carry
  at the penultimate position admits a closed-form expression
  \(c_{D-2} = \lfloor Q / 2^{2K-3} \rfloor - a_1 - c_1 - 2\), collapsing the
  \(O(K^2)\) schoolbook carry chain into \(O(1)\) (follows from the schoolbook
  carry recurrence). The limit matrix \(S_\infty\) as \(K \to \infty\) is
  analytically computable: its entries are elementary functions of
  \(\ln 2, \ln 3, \ln 5, \ln 7\) and rationals, and its eigenvalue constants are
  new transcendentals (these constants are not expressible as rational
  combinations of standard transcendentals, verified by PSLQ search). The
  complex structure is specific to the carry at position \(D-2\) --- carries at
  other positions yield purely real spectra. Unlike the Witt carry at odd primes
  (point 4), the D-odd binary carry eigenvectors are not orthogonal to
  \(\lvert 1\rangle\): the Artin-Schreier orthogonality mechanism is specific to
  \(p \geq 3\). The convergence
\end{enumerate}

\[R_{\mathrm{F}}(K) \to R_\infty \approx -0.341\]

has geometric rate exactly \(1/2\) (binary resolution doubling), distinct from
the eigenvalue modulus \(\lvert\lambda\rvert \approx 0.454\) governing Markov
decorrelation. (Here \(R_{\mathrm{F}}(K)\) denotes the D-odd carry ratio of the
4×4 top-bit matrix --- not to be confused with the cascade sector ratio
\(R(K) \to -\pi\) (Conjecture 1) of {[}P1, E{]}, which is a sum over all cascade
depths.)

Two obstructions prevent immediate globalization: the genus obstruction (the
\(p^2 \times p^2\) matrix captures only a subfactor of the degree-\(2g\)
Frobenius, where \(g = (p-1)(p-2)/2\)) and the cross-prime inconsistency (no
single Dirichlet character produces a coherent \(L\)-function across primes ---
different primes require different character orders). The Witt carry at
\(p = 2\) is degenerate (genus 0, carry always zero), so the connection between
the Archimedean trace anomaly (\(-\pi\)) and the non-Archimedean Frobenius
factors does not pass through the level-2 Witt carry.

\textbf{Open problems.}

\begin{itemize}
\tightlist
\item
  Identify the arithmetic rule that determines which character order is
  productive at each prime \(p\).
\item
  Compute the cohomology of the Artin-Schreier surface
  \(z^p - z = c_2(x_0, x_1, y_0, y_1)\) and verify that its Frobenius matches
  the carry spectrum at \(p = 3\) and \(p = 7\).
\item
  Explain the origin of the complex eigenvalues
  (\(\lambda \approx -0.302 \pm 0.340\,i\)) in the D-odd binary carry matrix.
  The limit matrix \(S_\infty\) is analytically known; the eigenvalues are new
  transcendental constants. The complex structure is specific to carry at
  \(D-2\) (other positions give real spectra). The relation to the oscillatory
  convergence \(R(K) \to -\pi\) of {[}E{]} remains open: the 4×4 matrix governs
  top-bit conditioning, while \(R(K)\) emerges from the full carry cascade.
\item
  Determine whether a family of Artin-Schreier varieties parametrized by \(p\)
  and a character can produce cross-prime-consistent Euler factors.
\end{itemize}

\begin{center}\rule{0.5\linewidth}{0.5pt}\end{center}

\newpage

\subsection{References}\label{references}

\begin{itemize}
\tightlist
\item
  {[}Del74{]} P. Deligne, \emph{La conjecture de Weil. I}, Publ. Math. IHÉS
  \textbf{43} (1974), 273--307.
\item
  {[}DF09{]} P. Diaconis and J. Fulman, \emph{Carries, shuffling, and symmetric
  functions}, Advances in Applied Mathematics \textbf{43} (2009), 176--196.
\item
  {[}Hol97{]} J. M. Holte, \emph{Carries, combinatorics, and an amazing matrix},
  American Mathematical Monthly \textbf{104} (1997), 138--149.
\item
  {[}IR90{]} K. Ireland and M. Rosen, \emph{A Classical Introduction to Modern
  Number Theory}, 2nd ed., Graduate Texts in Mathematics 84, Springer, 1990.
\item
  {[}Rue04{]} D. Ruelle, \emph{Thermodynamic Formalism}, 2nd ed., Cambridge
  University Press, 2004.
\item
  {[}Ser79{]} J.-P. Serre, \emph{Local Fields}, Graduate Texts in Mathematics
  67, Springer, 1979.
\item
  {[}Sil09{]} J. H. Silverman, \emph{The Arithmetic of Elliptic Curves}, 2nd
  ed., Graduate Texts in Mathematics 106, Springer, 2009.
\item
  {[}Wei48{]} A. Weil, \emph{Sur les courbes algébriques et les variétés qui
  s'en déduisent}, Actualités Sci. Ind. 1041, Hermann, 1948.
\item
  {[}Wei49{]} A. Weil, \emph{Numbers of solutions of equations in finite
  fields}, Bull. Amer. Math. Soc. \textbf{55} (1949), 497--508.
\item
  {[}Wit37{]} E. Witt, Zyklische Körper und Algebren der Characteristik \(p\)
  vom Grad \(p^n\), J. reine angew. Math. \textbf{176} (1937), 126--140.
\end{itemize}

\subsubsection{Companion Papers (carry-arithmetic
series)}\label{companion-papers-carry-arithmetic-series}

\begin{itemize}
\tightlist
\item
  {[}A{]} S. Alimonti, \emph{Spectral Theory of Carries in Positional
  Multiplication}, this series. (Foundation: the m-bit Equidistribution Lemma
  extending Diaconis--Fulman to the carry transfer operator.)
\item
  {[}E{]} S. Alimonti, \emph{The Trace Anomaly of Binary Multiplication}, this
  series. (The Shifted Resolvent Theorem: \(R = -\pi\) via resolvent
  universality, conditional on the Linear Mix Hypothesis.)
\item
  {[}G{]} S. Alimonti, \emph{The Angular Uniqueness of Base 2 in Positional
  Multiplication}, this series. (Why base 2 is the unique base for which
  \(c_1 = \pi/18\) involves \(\pi\); base-3 evidence for
  \(c_1(3) = \ln 3 - 1/2\).)
\item
  {[}P1{]} S. Alimonti, \emph{π from Pure Arithmetic: A Spectral Phase
  Transition in the Binary Carry Bridge}, this series. (The conjectured
  emergence of \(\pi\) from binary carry dynamics.)
\end{itemize}

\begin{center}\rule{0.5\linewidth}{0.5pt}\end{center}

\newpage

\subsection{Appendix A: Computational
Verification}\label{appendix-a-computational-verification}

All computations were performed in Python using exact arithmetic (the
\texttt{fractions.Fraction} type for \(\mathbb{Q}\)-coefficients and the
representation \(a + b\omega\) for \(\mathbb{Z}[\omega]\)). The
Faddeev--LeVerrier algorithm was used for the characteristic polynomial to avoid
the symbolic explosion of the determinant. Factorization was verified
independently using the computer algebra system SymPy.

The complete source code is available at
\texttt{experiments/F01\_verify\_theorem\_a.py}.

\subsubsection{A.1 The Exponent Matrix}\label{a.1-the-exponent-matrix}

The \(9 \times 9\) matrix \(E\) (with \(M_{j,i} = \omega^{E_{j,i}}\)) for
\(p = 3\), multiplicative Witt carry:

States ordered as
\((0,0), (0,1), (0,2), (1,0), (1,1), (1,2), (2,0), (2,1), (2,2)\).

\begin{verbatim}
E = [[0, 0, 0, 0, 0, 0, 0, 0, 0],
     [0, 0, 0, 0, 0, 0, 0, 0, 0],
     [0, 0, 0, 0, 0, 0, 0, 0, 0],
     [0, 0, 0, 0, 0, 0, 0, 0, 0],
     [1, 1, 1, 1, 2, 1, 1, 1, 0],
     [2, 2, 2, 2, 2, 1, 2, 0, 2],
     [0, 0, 0, 0, 0, 0, 0, 0, 0],
     [2, 2, 2, 2, 2, 0, 2, 1, 2],
     [1, 1, 1, 1, 0, 1, 1, 1, 2]]
\end{verbatim}

Entry distribution: 49 zeros, 16 ones, 16 twos.

\subsubsection{A.2 Reproduction}\label{a.2-reproduction}

\begin{Shaded}
\begin{Highlighting}[]
\ExtensionTok{pip}\NormalTok{ install numpy sympy}
\ExtensionTok{python}\NormalTok{ experiments/F01\_verify\_theorem\_a.py}
\end{Highlighting}
\end{Shaded}


\end{document}
